\iftestbuild % TEST-ONLY-BEGIN
\setcounter{section}{4}\setcounter{theorem}{86}
\fi % TEST-ONLY-END
% Source: book pp. 96--104 (PDF 110--118); solutions pp. 104--112, 172.
\section{Products and Quotients of Vector Spaces}

\subsection{Products of Vector Spaces}

As usual when dealing with more than one vector space, all vector spaces in use
should be over the same field.

\begin{definition}[Product of vector spaces]\label{ax:3.87}\leavevmode
Suppose $V_1,\dots,V_m$ are vector spaces over $\F$.
\begin{itemize}
\item The \emph{product} $V_1\times\cdots\times V_m$ is defined by
  \[ V_1\times\cdots\times V_m
     = \{(v_1,\dots,v_m) : v_1\in V_1,\dots,v_m\in V_m\}. \]
\item Addition on $V_1\times\cdots\times V_m$ is defined by
  \[ (u_1,\dots,u_m)+(v_1,\dots,v_m) = (u_1+v_1,\dots,u_m+v_m). \]
\item Scalar multiplication on $V_1\times\cdots\times V_m$ is defined by
  \[ \lambda(v_1,\dots,v_m) = (\lambda v_1,\dots,\lambda v_m). \]
\end{itemize}
\end{definition}

\begin{example}[Product of the vector spaces $\Poly_5(\R)$ and $\R^3$]\label{ax:3.88}
Elements of $\Poly_5(\R)\times\R^3$ are lists of length two, with the first item
in the list an element of $\Poly_5(\R)$ and the second item in the list an
element of $\R^3$.

For example, $\bigl(5-6x+4x^2,(3,8,7)\bigr)$ and $\bigl(x+9x^5,(2,2,2)\bigr)$
are elements of $\Poly_5(\R)\times\R^3$. Their sum is defined by
\begin{align*}
&\bigl(5-6x+4x^2,(3,8,7)\bigr)+\bigl(x+9x^5,(2,2,2)\bigr)\\
&\qquad= \bigl(5-5x+4x^2+9x^5,(5,10,9)\bigr).
\end{align*}
Also, $2\bigl(5-6x+4x^2,(3,8,7)\bigr)=\bigl(10-12x+8x^2,(6,16,14)\bigr)$.
\end{example}

The next result should be interpreted to mean that the product of vector spaces
is a vector space with the operations of addition and scalar multiplication as
defined by \ref{ax:3.87}.

\begin{result}[Product of vector spaces is a vector space]\label{ax:3.89}
Suppose $V_1,\dots,V_m$ are vector spaces over $\F$. Then
$V_1\times\cdots\times V_m$ is a vector space over $\F$.
\end{result}

The proof of the result above is left to the reader. Note that the additive
identity of $V_1\times\cdots\times V_m$ is $(0,\dots,0)$, where the $0$ in the
$k^{\text{th}}$ slot is the additive identity of $V_k$. The additive inverse of
$(v_1,\dots,v_m)\in V_1\times\cdots\times V_m$ is $(-v_1,\dots,-v_m)$.

\begin{example}[$\R^2\times\R^3\ne\R^5$ but $\R^2\times\R^3$ is isomorphic to
$\R^5$]\label{ax:3.90}
Elements of the vector space $\R^2\times\R^3$ are lists
\[ \bigl((x_1,x_2),(x_3,x_4,x_5)\bigr), \]
where $x_1,x_2,x_3,x_4,x_5\in\R$. Elements of $\R^5$ are lists
\[ (x_1,x_2,x_3,x_4,x_5), \]
where $x_1,x_2,x_3,x_4,x_5\in\R$.

Although elements of $\R^2\times\R^3$ and $\R^5$ look similar, they are not the
same kind of object. Elements of $\R^2\times\R^3$ are lists of length two (with
the first item itself a list of length two and the second item a list of length
three), and elements of $\R^5$ are lists of length five. Thus $\R^2\times\R^3$
does not equal $\R^5$.

The linear map
\[ \bigl((x_1,x_2),(x_3,x_4,x_5)\bigr) \mapsto (x_1,x_2,x_3,x_4,x_5) \]
is an isomorphism of the vector space $\R^2\times\R^3$ onto the vector space
$\R^5$. Thus these two vector spaces are isomorphic, although they are not
equal.
\end{example}

\marginnote[-30mm]{This isomorphism is so natural that we should think of it as a
relabeling. Some people informally say that $\R^2\times\R^3$ equals $\R^5$,
which is not technically correct but which captures the spirit of
identification via relabeling.}%
The next example illustrates the idea that we will use in the proof of
\ref{ax:3.92}.

\begin{example}[A basis of $\Poly_2(\R)\times\R^2$]\label{ax:3.91}
Consider this list of length five of elements of $\Poly_2(\R)\times\R^2$:
\[ \bigl(1,(0,0)\bigr),\ \bigl(x,(0,0)\bigr),\ \bigl(x^2,(0,0)\bigr),\
   \bigl(0,(1,0)\bigr),\ \bigl(0,(0,1)\bigr). \]
The list above is linearly independent and it spans $\Poly_2(\R)\times\R^2$.
Thus it is a basis of $\Poly_2(\R)\times\R^2$.
\end{example}

\begin{result}[Dimension of a product is the sum of dimensions]\label{ax:3.92}
Suppose $V_1,\dots,V_m$ are finite-dimensional vector spaces. Then
$V_1\times\cdots\times V_m$ is finite-dimensional and
\[ \dim(V_1\times\cdots\times V_m) = \dim V_1+\cdots+\dim V_m. \]
\end{result}

\begin{proof}
Choose a basis of each $V_k$. For each basis vector of each $V_k$, consider the
element of $V_1\times\cdots\times V_m$ that equals the basis vector in the
$k^{\text{th}}$ slot and $0$ in the other slots. The list of all such vectors is
linearly independent and spans $V_1\times\cdots\times V_m$. Thus it is a basis
of $V_1\times\cdots\times V_m$. The length of this basis is
$\dim V_1+\cdots+\dim V_m$, as desired.
\end{proof}

In the next result, the map $\Gamma$ is surjective by the definition of
$V_1+\cdots+V_m$. Thus the last word in the result below could be changed from
``injective'' to ``invertible''.

\begin{result}[Products and direct sums]\label{ax:3.93}
Suppose that $V_1,\dots,V_m$ are subspaces of $V$. Define a linear map
$\Gamma\colon V_1\times\cdots\times V_m\to V_1+\cdots+V_m$ by
\[ \Gamma(v_1,\dots,v_m) = v_1+\cdots+v_m. \]
Then $V_1+\cdots+V_m$ is a direct sum if and only if $\Gamma$ is injective.
\end{result}

\begin{proof}
By \ref{ax:3.15}, $\Gamma$ is injective if and only if the only way to write $0$
as a sum $v_1+\cdots+v_m$, where each $v_k$ is in $V_k$, is by taking each $v_k$
equal to $0$. Thus \ref{ax:1.45} shows that $\Gamma$ is injective if and only if
$V_1+\cdots+V_m$ is a direct sum, as desired.
\end{proof}

\begin{result}[A sum is a direct sum if and only if dimensions add
up]\label{ax:3.94}
Suppose $V$ is finite-dimensional and $V_1,\dots,V_m$ are subspaces of $V$. Then
$V_1+\cdots+V_m$ is a direct sum if and only if
\[ \dim(V_1+\cdots+V_m) = \dim V_1+\cdots+\dim V_m. \]
\end{result}

\begin{proof}
The map $\Gamma$ in \ref{ax:3.93} is surjective. Thus by the fundamental theorem
of linear maps (\ref{ax:3.21}), $\Gamma$ is injective if and only if
\[ \dim(V_1+\cdots+V_m) = \dim(V_1\times\cdots\times V_m). \]
Combining \ref{ax:3.93} and \ref{ax:3.92} now shows that $V_1+\cdots+V_m$ is a
direct sum if and only if
\[ \dim(V_1+\cdots+V_m) = \dim V_1+\cdots+\dim V_m, \]
as desired.
\end{proof}

In the special case $m=2$, an alternative proof that $V_1+V_2$ is a direct sum
if and only if $\dim(V_1+V_2)=\dim V_1+\dim V_2$ can be obtained by combining
\ref{ax:1.46} and \ref{ax:2.43}.

\subsection{Quotient Spaces}

We begin our approach to quotient spaces by defining the sum of a vector and a
subset.

\begin{notation}[$v+U$]\label{ax:3.95}
Suppose $v\in V$ and $U\subseteq V$. Then $v+U$ is the subset of $V$ defined by
\[ v+U = \{v+u : u\in U\}. \]
\end{notation}

\marginfig[16mm]{\begin{tikzpicture}[scale=.14]
  \draw[axis] (-1,0)--(21,0); \draw[axis] (0,-1.5)--(0,22);
  \draw[line width=.5pt] (-.6,-1.2)--(10.5,21) node[above,font=\scriptsize]
    {$(10,20)$};
  \draw[line width=.5pt] (6.4,-1.2)--(17.5,21) node[above right=-2pt,
    font=\scriptsize] {$(17,20)$};
  \draw (10,0)--(10,-.6) node[below,font=\scriptsize] {$10$};
  \draw (17,0)--(17,-.6) node[below,font=\scriptsize] {$17$};
  \draw (0,20)--(-.6,20) node[left,font=\scriptsize] {$20$};
  \node[font=\scriptsize] at (6.3,8.5) {$U$};
  \node[font=\scriptsize,right] at (13.6,11) {$(17,20)+U$};
\end{tikzpicture}}{$(17,20)+U$ is parallel to the subspace $U$.}%
\begin{example}[Sum of a vector and a one-dimensional subspace of
$\R^2$]\label{ax:3.96}
Suppose
\[ U = \{(x,2x)\in\R^2 : x\in\R\}. \]
Hence $U$ is the line in $\R^2$ through the origin with slope $2$. Thus
\[ (17,20)+U \]
is the line in $\R^2$ that contains the point $(17,20)$ and has slope $2$.

Because
\[ (10,20)\in U \quad\text{and}\quad (17,20)\in(17,20)+U, \]
we see that $(17,20)+U$ is obtained by moving $U$ to the right by $7$ units.
\end{example}

\begin{definition}[Translate]\label{ax:3.97}
For $v\in V$ and $U$ a subset of $V$, the set $v+U$ is said to be a
\emph{translate} of $U$.
\end{definition}

\begin{example}[Translates]\label{ax:3.98}\leavevmode
\begin{itemize}
\item If $U$ is the line in $\R^2$ defined by $U=\{(x,2x)\in\R^2 : x\in\R\}$,
  then all lines in $\R^2$ with slope $2$ are translates of $U$. See Example
  \ref{ax:3.96} above for a drawing of $U$ and one of its translates.
\item More generally, if $U$ is a line in $\R^2$, then the set of all
  translates of $U$ is the set of all lines in $\R^2$ that are parallel to $U$.
\item If $U=\{(x,y,0)\in\R^3 : x,y\in\R\}$, then the translates of $U$ are the
  planes in $\R^3$ that are parallel to the $xy$-plane $U$.
\item More generally, if $U$ is a plane in $\R^3$, then the set of all
  translates of $U$ is the set of all planes in $\R^3$ that are parallel to $U$
  (see, for example, Exercise~7).
\end{itemize}
\end{example}

\begin{definition}[Quotient space, $V/U$]\label{ax:3.99}
Suppose $U$ is a subspace of $V$. Then the \emph{quotient space} $V/U$ is the
set of all translates of $U$. Thus
\[ V/U = \{v+U : v\in V\}. \]
\end{definition}

\begin{example}[Quotient spaces]\label{ax:3.100}\leavevmode
\begin{itemize}
\item If $U=\{(x,2x)\in\R^2 : x\in\R\}$, then $\R^2/U$ is the set of all lines
  in $\R^2$ that have slope $2$.
\item If $U$ is a line in $\R^3$ containing the origin, then $\R^3/U$ is the set
  of all lines in $\R^3$ parallel to $U$.
\item If $U$ is a plane in $\R^3$ containing the origin, then $\R^3/U$ is the
  set of all planes in $\R^3$ parallel to $U$.
\end{itemize}
\end{example}

Our next goal is to make $V/U$ into a vector space. To do this, we will need
the next result.

\begin{result}[Two translates of a subspace are equal or
disjoint]\label{ax:3.101}
Suppose $U$ is a subspace of $V$ and $v,w\in V$. Then
\[ v-w\in U \iff v+U=w+U \iff (v+U)\cap(w+U)\ne\varnothing. \]
\end{result}

\begin{proof}
First suppose $v-w\in U$. If $u\in U$, then
\[ v+u = w+\bigl((v-w)+u\bigr)\in w+U. \]
Thus $v+U\subseteq w+U$. Similarly, $w+U\subseteq v+U$. Thus $v+U=w+U$,
completing the proof that $v-w\in U$ implies $v+U=w+U$.

The equation $v+U=w+U$ implies that $(v+U)\cap(w+U)\ne\varnothing$.

Now suppose $(v+U)\cap(w+U)\ne\varnothing$. Thus there exist $u_1,u_2\in U$
such that
\[ v+u_1 = w+u_2. \]
Thus $v-w=u_2-u_1$. Hence $v-w\in U$, showing that
$(v+U)\cap(w+U)\ne\varnothing$ implies $v-w\in U$, which completes the proof.
\end{proof}

Now we can define addition and scalar multiplication on $V/U$.

\begin{definition}[Addition and scalar multiplication on $V/U$]\label{ax:3.102}
Suppose $U$ is a subspace of $V$. Then \emph{addition} and \emph{scalar
multiplication} are defined on $V/U$ by
\begin{align*}
(v+U)+(w+U) &= (v+w)+U\\
\lambda(v+U) &= (\lambda v)+U
\end{align*}
for all $v,w\in V$ and all $\lambda\in\F$.
\end{definition}

As part of the proof of the next result, we will show that the definitions above
make sense.

\begin{result}[Quotient space is a vector space]\label{ax:3.103}
Suppose $U$ is a subspace of $V$. Then $V/U$, with the operations of addition
and scalar multiplication as defined above, is a vector space.
\end{result}

\begin{proof}
The potential problem with the definitions above of addition and scalar
multiplication on $V/U$ is that the representation of a translate of $U$ is not
unique. Specifically, suppose $v_1,v_2,w_1,w_2\in V$ are such that
\[ v_1+U=v_2+U \quad\text{and}\quad w_1+U=w_2+U. \]
To show that the definition of addition on $V/U$ given above makes sense, we
must show that $(v_1+w_1)+U=(v_2+w_2)+U$.

By \ref{ax:3.101}, we have
\[ v_1-v_2\in U \quad\text{and}\quad w_1-w_2\in U. \]
Because $U$ is a subspace of $V$ and thus is closed under addition, this implies
that $(v_1-v_2)+(w_1-w_2)\in U$. Thus $(v_1+w_1)-(v_2+w_2)\in U$. Using
\ref{ax:3.101} again, we see that
\[ (v_1+w_1)+U = (v_2+w_2)+U, \]
as desired. Thus the definition of addition on $V/U$ makes sense.

Similarly, suppose $\lambda\in\F$. We are still assuming that $v_1+U=v_2+U$.
Because $U$ is a subspace of $V$ and thus is closed under scalar multiplication,
we have $\lambda(v_1-v_2)\in U$. Thus $\lambda v_1-\lambda v_2\in U$. Hence
\ref{ax:3.101} implies that
\[ (\lambda v_1)+U = (\lambda v_2)+U. \]
Thus the definition of scalar multiplication on $V/U$ makes sense.

Now that addition and scalar multiplication have been defined on $V/U$, the
verification that these operations make $V/U$ into a vector space is
straightforward and is left to the reader. Note that the additive identity of
$V/U$ is $0+U$ (which equals $U$) and that the additive inverse of $v+U$ is
$(-v)+U$.
\end{proof}

The next concept will lead to a computation of the dimension of $V/U$.

\begin{definition}[Quotient map, $\pi$]\label{ax:3.104}
Suppose $U$ is a subspace of $V$. The \emph{quotient map} $\pi\colon V\to V/U$
is the linear map defined by
\[ \pi(v) = v+U \]
for each $v\in V$.
\end{definition}

The reader should verify that $\pi$ is indeed a linear map. Although $\pi$
depends on $U$ as well as $V$, these spaces are left out of the notation because
they should be clear from the context.

\begin{result}[Dimension of quotient space]\label{ax:3.105}
Suppose $V$ is finite-dimensional and $U$ is a subspace of $V$. Then
\[ \dim V/U = \dim V-\dim U. \]
\end{result}

\begin{proof}
Let $\pi$ denote the quotient map from $V$ to $V/U$. If $v\in V$, then
$v+U=0+U$ if and only if $v\in U$ (by \ref{ax:3.101}), which implies that
$\nul\pi=U$. The definition of $\pi$ implies $\range\pi=V/U$. The fundamental
theorem of linear maps (\ref{ax:3.21}) now implies $\dim V=\dim U+\dim V/U$,
which gives the desired result.
\end{proof}

Each linear map $T$ on $V$ induces a linear map $\widetilde{T}$ on
$V/(\nul T)$, as defined below.

\begin{notation}[$\widetilde{T}$]\label{ax:3.106}
Suppose $T\in\Lin(V,W)$. Define $\widetilde{T}\colon V/(\nul T)\to W$ by
\[ \widetilde{T}(v+\nul T) = Tv. \]
\end{notation}

To show that the definition of $\widetilde{T}$ makes sense, suppose $u,v\in V$
are such that $u+\nul T=v+\nul T$. By \ref{ax:3.101}, we have $u-v\in\nul T$.
Thus $T(u-v)=0$. Hence $Tu=Tv$. Thus the definition of $\widetilde{T}$ indeed
makes sense. The routine verification that $\widetilde{T}$ is a linear map from
$V/(\nul T)$ to $W$ is left to the reader.

The next result shows that we can think of $\widetilde{T}$ as a modified version
of $T$, with a domain that produces a one-to-one map.

\begin{result}[Null space and range of $\widetilde{T}$]\label{ax:3.107}
Suppose $T\in\Lin(V,W)$. Then
\begin{parts}
\item $\widetilde{T}\circ\pi=T$, where $\pi$ is the quotient map of $V$ onto
  $V/(\nul T)$;
\item $\widetilde{T}$ is injective;
\item $\range\widetilde{T}=\range T$;
\item $V/(\nul T)$ and $\range T$ are isomorphic vector spaces.
\end{parts}
\end{result}

\begin{proof}\leavevmode
\begin{parts}
\item If $v\in V$, then $(\widetilde{T}\circ\pi)(v)=\widetilde{T}\bigl(\pi(v)\bigr)
  =\widetilde{T}(v+\nul T)=Tv$, as desired.
\item Suppose $v\in V$ and $\widetilde{T}(v+\nul T)=0$. Then $Tv=0$. Thus
  $v\in\nul T$. Hence \ref{ax:3.101} implies that $v+\nul T=0+\nul T$. This
  implies that $\nul\widetilde{T}=\{0+\nul T\}$. Hence $\widetilde{T}$ is
  injective, as desired.
\item The definition of $\widetilde{T}$ shows that
  $\range\widetilde{T}=\range T$.
\item Now (b) and (c) imply that if we think of $\widetilde{T}$ as mapping into
  $\range T$, then $\widetilde{T}$ is an isomorphism from $V/(\nul T)$ onto
  $\range T$. \qedhere
\end{parts}
\end{proof}

% ------------------------------------------------------------------ exercises
\exercisesheading{3E}

\begin{exercise}
Suppose $T$ is a function from $V$ to $W$. The \emph{graph} of $T$ is the
subset of $V\times W$ defined by
\[ \text{graph of } T = \{(v,Tv)\in V\times W : v\in V\}. \]
Prove that $T$ is a linear map if and only if the graph of $T$ is a subspace of
$V\times W$.
\exnote{Formally, a function $T$ from $V$ to $W$ is a subset $T$ of $V\times W$
such that for each $v\in V$, there exists exactly one element $(v,w)\in T$. In
other words, formally a function is what is called above its graph. We do not
usually think of functions in this formal manner. However, if we do become
formal, then this exercise could be rephrased as follows: Prove that a function
$T$ from $V$ to $W$ is a linear map if and only if $T$ is a subspace of
$V\times W$.}
\end{exercise}
\begin{solution}\emergencystretch=1.5em
Let $G(T)$ denote the graph of $T$. Now, suppose $T$ is a linear map. Then
$G(T)$ is nonempty since $(0,T0)=(0,0)\in G(T)$. Now, let
$(v_1,Tv_1),(v_2,Tv_2)\in G(T)$. Then
\[ (v_1-v_2,Tv_1-Tv_2) = (v_1,Tv_1)-(v_2,Tv_2)\in G(T) \]
since $T$ is linear. Now, let $(v,Tv)\in G(T)$ and $\alpha\in\F$. Then
\[ (\alpha v,\alpha Tv) = \alpha(v,Tv)\in G(T) \]
since $T$ is linear. Thus, $G(T)$ is a subspace of $V\times W$.

Conversely, suppose $G(T)$ is a subspace of $V\times W$. By definition, we have
$(v_1,Tv_1),\allowbreak(v_2,Tv_2)\in G(T)$ for all $v_1,v_2\in V$. Since $G(T)$ is a
subspace of $V\times W$, for all $\alpha,\beta\in\F$, we have
\[ (\alpha v_1+\beta v_2,\alpha Tv_1+\beta Tv_2)
   = \alpha(v_1,Tv_1)+\beta(v_2,Tv_2)\in G(T). \]
By definition of $G(T)$, we have $\alpha Tv_1+\beta Tv_2=T(\alpha v_1+\beta
v_2)$. Thus, $T$ is a linear map.
\end{solution}

\begin{exercise}
Suppose that $V_1,\dots,V_m$ are vector spaces such that
$V_1\times\cdots\times V_m$ is finite-dimensional. Prove that $V_k$ is
finite-dimensional for each $k=1,\dots,m$.
\end{exercise}
\begin{solution}\emergencystretch=1.5em
Let $V=V_1\times\cdots\times V_m$, and consider $T\in\Lin(V,V_k)$ defined by
$T(v_1,\dots,v_m)=v_k$. Then $T$ is a linear map since for all
$\alpha,\beta\in\F$ and all $(v_1,\dots,v_m),(w_1,\dots,w_m)\in V$, we have
\begin{align*}
T\bigl(\alpha(v_1,\dots,v_m)+\beta(w_1,\dots,w_m)\bigr)
  &= T(\alpha v_1+\beta w_1,\dots,\alpha v_m+\beta w_m)\\
  &= \alpha v_k+\beta w_k\\
  &= \alpha T(v_1,\dots,v_m)+\beta T(w_1,\dots,w_m).
\end{align*}
Moreover, $T$ is surjective since for any $v_k\in V_k$, we have
$v_k=T(0,\dots,0,v_k,0,\dots,0)$ where $v_k$ is in the $k$-th coordinate. Since
$V$ is finite-dimensional and $T$ is a linear map that is surjective, it follows
that $V_k$ is finite-dimensional.
\end{solution}

\begin{exercise}
Suppose $V_1,\dots,V_m$ are vector spaces. Prove that
$\Lin(V_1\times\cdots\times V_m,W)$ and
$\Lin(V_1,W)\times\cdots\times\Lin(V_m,W)$ are isomorphic vector spaces.
\exnote{There is no assumption in the exercise above or in the two following
exercises that the vector spaces are finite-dimensional.}
\end{exercise}
\begin{solution}
Let $V=V_1\times\cdots\times V_m$. Let
$L=\Lin(V_1,W)\times\cdots\times\Lin(V_m,W)$. Define
$T\in\Lin\bigl(\Lin(V,W),L\bigr)$ by
\[ T(S) = (S\circ\iota_1,\dots,S\circ\iota_m) \]
where $\iota_k\in\Lin(V_k,V)$ is the map defined as
\[ \iota_k(v_k) = (0,\dots,0,v_k,0,\dots,0) \]
with $\iota_k(v_k)$ having $v_k$ in the $k$-th coordinate and $0$ in all other
coordinates. Then $T$ is linear since for all $\alpha,\beta\in\F$ and all
$S_1,S_2\in\Lin(V,W)$, we have
\begin{align*}
T(\alpha S_1+\beta S_2)
  &= \bigl((\alpha S_1+\beta S_2)\circ\iota_1,\dots,
     (\alpha S_1+\beta S_2)\circ\iota_m\bigr)\\
  &= \bigl(\alpha(S_1\circ\iota_1)+\beta(S_2\circ\iota_1),\dots,
     \alpha(S_1\circ\iota_m)+\beta(S_2\circ\iota_m)\bigr)\\
  &= \alpha T(S_1)+\beta T(S_2),
\end{align*}
hence $T$ is a linear map. Now suppose $T(S)=0$. Then $S\circ\iota_k=0$ for all
$k=1,\dots,m$. Moreover, for all $(v_1,\dots,v_m)\in V$, we have
\begin{align*}
S(v_1,\dots,v_m) &= S\bigl(\iota_1(v_1)+\cdots+\iota_m(v_m)\bigr)\\
  &= S\circ\iota_1(v_1)+\cdots+S\circ\iota_m(v_m) = 0.
\end{align*}
It follows that $S(v_1,\dots,v_m)=0$ for all $(v_1,\dots,v_m)\in V$. Thus,
$S=0$, and hence $T$ is injective.

Now, let $(S_1,\dots,S_m)\in L$. Define $S\in\Lin(V,W)$ by
$S(v_1,\dots,v_m)=S_1v_1+\cdots+S_mv_m$. Then $T(S)=(S_1,\dots,S_m)$. Thus, $T$
is surjective. Therefore, $T$ is an isomorphism.
\end{solution}

\begin{exercise}
Suppose $W_1,\dots,W_m$ are vector spaces. Prove that
$\Lin(V,W_1\times\cdots\times W_m)$ and
$\Lin(V,W_1)\times\cdots\times\Lin(V,W_m)$ are isomorphic vector spaces.
\end{exercise}
\begin{solution}\emergencystretch=1.5em
\solnote{Manual Remark}{Note the difference between this exercise and
Exercise~3E3. Here, the domain is decomposed, while in Exercise~3E3, the
codomain is decomposed. Hence, the use of the maps $\iota_k$ and $\pi_k$ in the
solutions to these exercises are reversed.}
Let $W=W_1\times\cdots\times W_m$. Let
$L=\Lin(V,W_1)\times\cdots\times\Lin(V,W_m)$. Define
$T\in\Lin\bigl(\Lin(V,W),L\bigr)$ by
\[ T(S) = (\pi_1\circ S,\dots,\pi_m\circ S) \]
where $\pi_k\in\Lin(W,W_k)$ is the map defined as
\[ \pi_k(w_1,\dots,w_m) = w_k \]
Then $T$ is linear since for all $\alpha,\beta\in\F$ and all
$S_1,S_2\in\Lin(V,W)$, we have
\begin{align*}
T(\alpha S_1+\beta S_2)
  &= \bigl(\pi_1\circ(\alpha S_1+\beta S_2),\dots,
     \pi_m\circ(\alpha S_1+\beta S_2)\bigr)\\
  &= \bigl(\alpha(\pi_1\circ S_1)+\beta(\pi_1\circ S_2),\dots,
     \alpha(\pi_m\circ S_1)+\beta(\pi_m\circ S_2)\bigr)\\
  &= \alpha T(S_1)+\beta T(S_2),
\end{align*}
hence $T$ is a linear map. Now suppose $T(S)=0$. Then $\pi_k\circ S=0$ for all
$k=1,\dots,m$. Moreover, for all $v\in V$, we have
\[ Sv = \bigl(\pi_1(Sv),\dots,\pi_m(Sv)\bigr) = (0,\dots,0) = 0. \]
It follows that $Sv=0$ for all $v$, hence $S=0$, and thus $T$ is injective.

Now let $(S_1,\dots,S_m)\in L$. Define $S\in\Lin(V,W)$ by
$Sv=(S_1v,\dots,S_mv)$. Then $T(S)=(S_1,\dots,S_m)$. Thus, $T$ is surjective.
Therefore, $T$ is an isomorphism.
\end{solution}

\begin{exercise}
For $m$ a positive integer, define $V^m$ by
\[ V^m = \ftimes{V\times\cdots\times V}{m\text{ times}}. \]
Prove that $V^m$ and $\Lin(\F^m,V)$ are isomorphic vector spaces.
\end{exercise}
\begin{solution}\emergencystretch=1.5em
Let $e_1,\dots,e_m$ be the standard basis of $\F^m$. Define
$T\in\Lin\bigl(V^m,\allowbreak\Lin(\F^m,V)\bigr)$ by
\[ T(v_1,\dots,v_m)(\lambda_1,\dots,\lambda_m)
   = \lambda_1v_1+\cdots+\lambda_mv_m. \]
Then $T$ is linear since for all $\alpha,\beta\in\F$ and all
$v=(v_1,\dots,v_m),w=(w_1,\dots,w_m)\in V^m$, we have
\begin{align*}
T(\alpha v+\beta w)(\lambda_1,\dots,\lambda_m)
  &= T(\alpha v_1+\beta w_1,\dots,\alpha v_m+\beta w_m)
     (\lambda_1,\dots,\lambda_m)\\
  &= \lambda_1(\alpha v_1+\beta w_1)+\cdots+\lambda_m(\alpha v_m+\beta w_m)\\
  &= \alpha(\lambda_1v_1+\cdots+\lambda_mv_m)
     +\beta(\lambda_1w_1+\cdots+\lambda_mw_m)\\
  &= \alpha T(v)(\lambda_1,\dots,\lambda_m)
     +\beta T(w)(\lambda_1,\dots,\lambda_m),
\end{align*}
hence $T$ is a linear map. Now suppose $T(v)=0$. Then $T(v)(e_k)=0$ for all
$k=1,\dots,m$. It follows that $v_k=0$ for all $k=1,\dots,m$, hence $v=0$, and
thus $T$ is injective.

Now let $S\in\Lin(\F^m,V)$. Define $v\in V^m$ by $v_k=Se_k$ for all
$k=1,\dots,m$. Then $T(v)=S$. Thus, $T$ is surjective. Therefore, $T$ is an
isomorphism.
\end{solution}

\begin{exercise}
Suppose that $v,x$ are vectors in $V$ and that $U,W$ are subspaces of $V$ such
that $v+U=x+W$. Prove that $U=W$.
\end{exercise}
\begin{solution}
Let $u\in U$. Then $v+u\in v+U=x+W$. Thus, there exists $w\in W$ such that
$v+u=x+w$. Note that $v\in x+W$ implies the existence of some $w'\in W$ such
that $v=x+w'$, or $x-v=-w'\in W$. Similarly, $x\in v+U$ implies the existence of
some $u'\in U$ such that $x=v+u'$, or $v-x=-u'\in U$. It follows that
$u=(x-v)+w\in W$, hence $U\subseteq W$. Now let $w\in W$. Then
$x+w\in x+W=v+U$. Thus, there exists $u\in U$ such that $x+w=v+u$. It follows
that $w=(v-x)+u\in U$, hence $W\subseteq U$. Therefore, $U=W$.
\end{solution}

\begin{exercise}
Let $U=\{(x,y,z)\in\R^3 : 2x+3y+5z=0\}$. Suppose $A\subseteq\R^3$. Prove that
$A$ is a translate of $U$ if and only if there exists $c\in\R$ such that
\[ A = \{(x,y,z)\in\R^3 : 2x+3y+5z=c\}. \]
\end{exercise}
\begin{solution}
Suppose $A$ is a translate of $U$. Then there exists $v\in\R^3$ such that
$A=v+U$. Let $v=(x_0,y_0,z_0)$. Then
\begin{align*}
A &= v+U\\
  &= \{(x_0,y_0,z_0)+(x,y,z)\in\R^3 \mid 2x+3y+5z=0\}\\
  &= \{(x',y',z')\in\R^3 \mid 2(x'-x_0)+3(y'-y_0)+5(z'-z_0)=0\}\\
  &= \{(x',y',z')\in\R^3 \mid 2x'+3y'+5z'=2x_0+3y_0+5z_0\}.
\end{align*}
Conversely, suppose there exists $c\in\R$ such that
\[ A = \{(x,y,z)\in\R^3 \mid 2x+3y+5z=c\}. \]
Let $v=(x_0,y_0,z_0)$ where $x_0,y_0,z_0$ are any real numbers such that
$2x_0+3y_0+5z_0=c$. Then
\begin{align*}
v+U &= (x_0,y_0,z_0)+U\\
    &= \{(x_0,y_0,z_0)+(x,y,z)\in\R^3 \mid 2x+3y+5z=0\}\\
    &= A. \qedhere
\end{align*}
\end{solution}

\begin{exercise}\leavevmode
\begin{parts}
\item Suppose $T\in\Lin(V,W)$ and $c\in W$. Prove that $\{x\in V : Tx=c\}$ is
  either the empty set or is a translate of $\nul T$.
\item Explain why the set of solutions to a system of linear equations such as
  \ref{ax:3.27} is either the empty set or is a translate of some subspace of
  $\F^n$.
\end{parts}
\end{exercise}
\begin{solution}\leavevmode
\begin{parts}
\item Suppose $T\in\Lin(V,W)$ and $c\in W$. Let $A=\{x\in V \mid Tx=c\}$. If $A$
  is nonempty, then there exists $v\in V$ such that $Tv=c$. We claim that
  $A=v+\nul T$. Let $x\in A$. Then $Tx=c=Tv$, hence $T(x-v)=0$, and thus
  $x-v\in\nul T$. It follows that $x\in v+\nul T$, hence $A\subseteq v+\nul T$.
  Now let $y\in v+\nul T$. Then there exists $u\in\nul T$ such that $y=v+u$. It
  follows that $Ty=T(v+u)=Tv+Tu=c+0=c$, hence $y\in A$, and thus
  $v+\nul T\subseteq A$. Therefore, $A=v+\nul T$.
\item A system of linear equations such as \ref{ax:3.27} can be written in the
  form $Tx=c$ for some linear map $T$ and some vector $c$. By part (a), the set
  of solutions to this system of linear equations is either the empty set or is
  a translate of $\nul T$, which is a subspace of $\F^n$. \qedhere
\end{parts}
\end{solution}

\begin{exercise}
Prove that a nonempty subset $A$ of $V$ is a translate of some subspace of $V$
if and only if $\lambda v+(1-\lambda)w\in A$ for all $v,w\in A$ and all
$\lambda\in\F$.
\end{exercise}
\begin{solution}
Suppose $A$ is a translate of some subspace of $V$. Then there exists $v\in V$
and a subspace $U$ of $V$ such that $A=v+U$. Let $v+u,v+u'\in A$. Then for any
$\lambda\in\F$,
\[ \lambda(v+u)+(1-\lambda)(v+u')
   = v+\underbrace{\bigl(\lambda u+(1-\lambda)u'\bigr)}_{\in U}\in A. \]
Conversely, suppose $\lambda v+(1-\lambda)w\in A$ for all $v,w\in A$ and all
$\lambda\in\F$. Fix some $v_0\in A$, and let $U=\{u\in V \mid v_0+u\in A\}$. We
first claim that $U$ is a subspace of $V$. Since $v_0=v_0+0\in A$, we have
$0\in U$. Let $u_1,u_2\in U$. Then $v_0+u_1,v_0+u_2\in A$. Now let
$\lambda\in\F$. Then
\begin{align*}
\lambda(v_0+u_1)+(1-\lambda)(v_0+u_2) &= v_0+\lambda u_1+(1-\lambda)u_2\\
  &= v_0+\bigl(\lambda u_1+(1-\lambda)u_2\bigr)\in A.
\end{align*}
It follows that $\lambda u_1+(1-\lambda)u_2\in U$, hence $U$ is a subspace of
$V$. Now we claim that $A=v_0+U$. Let $w\in A$. By definition,
$v_0+(w-v_0)=w\in A$, hence $w-v_0\in U$, and thus $w\in v_0+U$, so
$A\subseteq v_0+U$. Now let $v_0+u\in v_0+U$. Then $v_0+u\in A$ by definition of
$U$, hence $v_0+U\subseteq A$. Therefore, $A=v_0+U$.
\end{solution}

\begin{exercise}
Suppose $A_1=v+U_1$ and $A_2=w+U_2$ for some $v,w\in V$ and some subspaces
$U_1,U_2$ of $V$. Prove that the intersection $A_1\cap A_2$ is either a
translate of some subspace of $V$ or is the empty set.
\end{exercise}
\begin{solution}
If $A_1\cap A_2$ is the empty set, then we are done. Suppose that $A_1\cap A_2$
is nonempty. Let $x,y\in A_1\cap A_2$ and $\lambda\in\F$. Then $x,y\in A_1$,
hence $\lambda x+(1-\lambda)y\in A_1$. Moreover, $x,y\in A_2$, hence
$\lambda x+(1-\lambda)y\in A_2$. It follows that
$\lambda x+(1-\lambda)y\in A_1\cap A_2$. Thus, by Exercise~3E9, $A_1\cap A_2$ is
a translate of some subspace of $V$.
\end{solution}

\begin{exercise}
Suppose $U=\{(x_1,x_2,\dots)\in\F^\infty : x_k\ne0 \text{ for only finitely many
} k\}$.
\begin{parts}
\item Show that $U$ is a subspace of $\F^\infty$.
\item Prove that $\F^\infty/U$ is infinite-dimensional.
\end{parts}
\end{exercise}
\begin{solution}\emergencystretch=1.5em
\solnote{Appendix}{See the alternate proof at the end of this solution for a
generalization of this treatment, so long as you are comfortable with the idea
$\N$ can be partitioned into infinitely many disjoint infinite sets.}
\begin{parts}
\item We claim that $U$ can be reformulated as the set of vectors
  $(x_1,x_2,\dots)\in\F^\infty$ such that there exists an index $N$ such that
  $x_k=0$ for all $k>N$. To see this, suppose $(x_1,x_2,\dots)\in U$. Then there
  exists a finite set of indices $K$ such that $x_k=0$ for all $k\notin K$. If
  $K$ is empty, let $N=0$. Otherwise, let $N=\max K$. Then $x_k=0$ for all
  $k>N$.

  Conversely, suppose there exists an index $N$ such that $x_k=0$ for all
  $k>N$. Then the set of indices $\{1,2,\dots,N\}$ is finite and contains all
  indices where $x_k\ne0$, hence $(x_1,x_2,\dots)\in U$. Thus, the two
  formulations of $U$ are equivalent.

  We now show that $U$ is a subspace of $\F^\infty$. First, note that the zero
  vector $(0,0,\dots)\in U$ since it has no nonzero coordinates. Now let
  $u=(u_1,u_2,\dots),v=(v_1,v_2,\dots)\in U$. Then there exist indices $N_u$ and
  $N_v$ such that $u_k=0$ for all $k>N_u$ and $v_k=0$ for all $k>N_v$. Let
  $N=\max(N_u,N_v)$. Then for all $k>N$, we have $u_k=0$ and $v_k=0$, hence
  $(u+v)_k=u_k+v_k=0+0=0$. Thus, $u+v\in U$. Now let $\alpha\in\F$. Then for all
  $k>N_u$, we have $(\alpha u)_k=\alpha u_k=\alpha\cdot0=0$, hence
  $\alpha u\in U$. Therefore, $U$ is a subspace of $\F^\infty$.
\item For any given positive integer $m$, define the sets $B_i$ for
  $1\le i\le m$ as follows:
  \[ B_i = \{i+km \mid k\in\N\} = \{i,i+m,i+2m,\dots\}. \]
  Then $B_1,\dots,B_m$ are pairwise disjoint subsets of $\N$. For each
  $i=1,\dots,m$, define
  \[ v_i = (v_{i,1},v_{i,2},\dots)\in\F^\infty \]
  such that $v_{i,j}=1$ if $j\in B_i$ and $v_{i,j}=0$ otherwise. We claim that
  $v_1+U,\dots,v_m+U$ are linearly independent in $\F^\infty/U$.

  Let $\lambda_1,\dots,\lambda_m\in\F$ such that
  $\lambda_1(v_1+U)+\cdots+\lambda_m(v_m+U)=U$. Then
  $\lambda_1v_1+\cdots+\lambda_mv_m\in U$. Then we can show that $\lambda_i=0$
  for all $i=1,\dots,m$ as follows. For each $i$, consider the $j$-th coordinate
  of $v_i$ for $j\in B_i$. Since $v_i$ has $1$'s in all coordinates indexed by
  $B_i$ and $0$'s elsewhere, and $B_i$ is disjoint from $B_k$ for $k\ne i$, the
  $j$-th coordinate of $\lambda_1v_1+\cdots+\lambda_mv_m$ is $\lambda_i$. But
  this sum is in $U$, which means it has only finitely many nonzero coordinates.
  Since $B_i$ is infinite, it follows that $\lambda_i=0$. Hence, all
  $\lambda_i=0$, proving that $v_1+U,\dots,v_m+U$ are linearly independent.
  Since $m$ was arbitrary, $\F^\infty/U$ is infinite-dimensional.
\end{parts}

\emph{Alternate proof.}
\begin{parts}
\item Since $(0,0,\dots)\in U$, $U$ is nonempty. Suppose $x,y\in U$ and
  $c\in\F$. Let $S_x$ and $S_y$ be the sets of indices where $x$ and $y$ have
  nonzero terms, respectively. By definition, $S_x$ and $S_y$ are finite.
  $cx+y$ can only have nonzero terms at indices in $S_x\cup S_y$. Since the
  union of two finite sets is finite, $cx+y\in U$, which makes $U$ a subspace
  of $\F^\infty$.
\item Partition $\Z^+$ into a countably infinite family of disjoint infinite
  subsets $A_1,\allowbreak A_2,\allowbreak\dots$. For each $n$, define $v_n\in\F^\infty$ such that its
  $j$-th term is $1$ if $j\in A_n$ and $0$ otherwise.

  Suppose that $\sum_{i=1}^m c_i(v_i+U)=0+U$ for some positive integer $m$ and
  scalars $c_1,\dots,c_m\in\F$. Then, $w=\sum_{i=1}^m c_iv_i\in U$. By
  construction, for $i\in\{1,\dots,m\}$, the $j$-th term of $w$ is $c_i$ for all
  $j\in A_i$. Since $w\in U$, it only has a finite number of nonzero terms. But
  each $A_i$ is infinite, so this forces $c_i=0$ for all $i\in\{1,\dots,m\}$.
  Hence, $v_1+U,v_2+U,\dots$ are linearly independent in $\F^\infty/U$.

  Because $\F^\infty/U$ contains an infinite sequence of linearly independent
  vectors, it is infinite-dimensional.
\end{parts}
\end{solution}

\begin{exercise}
Suppose $v_1,\dots,v_m\in V$. Let
\[ A = \{\lambda_1v_1+\cdots+\lambda_mv_m : \lambda_1,\dots,\lambda_m\in\F
   \text{ and } \lambda_1+\cdots+\lambda_m=1\}. \]
\begin{parts}
\item Prove that $A$ is a translate of some subspace of $V$.
\item Prove that if $B$ is a translate of some subspace of $V$ and
  $\{v_1,\dots,v_m\}\subseteq B$, then $A\subseteq B$.
\item Prove that $A$ is a translate of some subspace of $V$ of dimension less
  than $m$.
\end{parts}
\end{exercise}
\begin{solution}\leavevmode
\begin{parts}
\item For any element $\lambda_1v_1+\cdots+\lambda_mv_m\in A$, we can write
  \begin{align*}
  \lambda_1v_1+\cdots+\lambda_mv_m
    &= v_1+\lambda_2(v_2-v_1)+\cdots+\lambda_m(v_m-v_1)\\
    &= v_1+\sum_{i=2}^m\lambda_i(v_i-v_1),
  \end{align*}
  where $\lambda_1=1-(\lambda_2+\cdots+\lambda_m)$. Let
  $U=\spn(v_2-v_1,\dots,v_m-v_1)$. Then $A=v_1+U$, hence $A$ is a translate of
  the subspace $U$ of $V$.
\item By Exercise~3E9, since $B$ is a translate of some subspace of $V$, we have
  $\lambda v+(1-\lambda)w\in B$ for all $v,w\in B$ and all $\lambda\in\F$. Since
  $\{v_1,\dots,v_m\}\subseteq B$, it follows that any linear combination of
  $v_1,\dots,v_m$ with coefficients summing to $1$ is in $B$. Hence,
  $A\subseteq B$.
\item The subspace $U=\spn(v_2-v_1,\dots,v_m-v_1)$ has dimension at most $m-1$,
  since it is spanned by $m-1$ vectors. Therefore, $A=v_1+U$ is a translate of a
  subspace of $V$ of dimension less than $m$. \qedhere
\end{parts}
\end{solution}

\begin{exercise}
Suppose $U$ is a subspace of $V$ such that $V/U$ is finite-dimensional. Prove
that $V$ is isomorphic to $U\times(V/U)$.
\end{exercise}
\begin{solution}
Since $V/U$ is finite-dimensional, let $v_1+U,\dots,v_m+U$ be a basis of $V/U$.
Let $W=\spn(v_1,\dots,v_m)$. We show that $V=U\oplus W$. First, note that any
$v\in V$ can be written as $v=u+w$ for some $u\in U$ and $w\in W$. Indeed, since
$v+U\in V/U$, we can write
\[ v+U = \lambda_1(v_1+U)+\cdots+\lambda_m(v_m+U) \]
for some scalars $\lambda_1,\dots,\lambda_m$. This gives
\[ v-(\lambda_1v_1+\cdots+\lambda_mv_m)\in U. \]
So setting $u=v-(\lambda_1v_1+\cdots+\lambda_mv_m)\in U$ and
$w=\lambda_1v_1+\cdots+\lambda_mv_m\in W$, we have $v=u+w$. To see that the sum
is direct, suppose $u+w=0$ with $u\in U$ and $w\in W$. Then $w=-u\in W\cap U$.
Since $w\in W$, we can write $w=\lambda_1v_1+\cdots+\lambda_mv_m$ for some
scalars $\lambda_1,\dots,\lambda_m$. Since we established that $w\in U$, we have
$w+U=0+U$. But $v_1+U,\dots,v_m+U$ are linearly independent in $V/U$, hence
$\lambda_1=\cdots=\lambda_m=0$, and thus $w=0$. It follows that $u=0$ as well,
proving that the sum is direct. Therefore, $V=U\oplus W$.

Consider the map $T\colon V\to U\times(V/U)$ given by $Tv=(u,v+U)$, where
$v=u+w$ with $u\in U$ and $w\in W$. This map is well-defined because the
decomposition $v=u+w$ is unique. It is linear, injective, and surjective, hence
an isomorphism. Therefore, $V\cong U\times(V/U)$.
\end{solution}

\begin{exercise}
Suppose $U$ and $W$ are subspaces of $V$ and $V=U\oplus W$. Suppose
$w_1,\dots,w_m$ is a basis of $W$. Prove that $w_1+U,\dots,w_m+U$ is a basis of
$V/U$.
\end{exercise}
\begin{solution}
Since $V=U\oplus W$, every $v\in V$ can be written uniquely as $v=u+w$ with
$u\in U$ and $w\in W$. We claim that $w_1+U,\dots,w_m+U$ is a basis of $V/U$.
First, we show that they span $V/U$. Let $v+U\in V/U$. Write $v=u+w$ with
$u\in U$ and $w\in W$. Then $v+U=(u+w)+U=w+U$. Since $w\in W$, we can write
$w=\lambda_1w_1+\cdots+\lambda_mw_m$ for some scalars
$\lambda_1,\dots,\lambda_m$. Hence,
$v+U=\lambda_1(w_1+U)+\cdots+\lambda_m(w_m+U)$, showing that
$w_1+U,\dots,w_m+U$ span $V/U$.

Next, we show linear independence. Suppose
$\lambda_1(w_1+U)+\cdots+\lambda_m(w_m+U)=0+U$. Then
$\lambda_1w_1+\cdots+\lambda_mw_m\in U$. But
$\lambda_1w_1+\cdots+\lambda_mw_m\in W$ and $U\cap W=\{0\}$, so
$\lambda_1w_1+\cdots+\lambda_mw_m=0$. Since $w_1,\dots,w_m$ are linearly
independent, we have $\lambda_1=\cdots=\lambda_m=0$. Therefore,
$w_1+U,\dots,w_m+U$ are linearly independent, and hence form a basis of $V/U$.
\end{solution}

\begin{exercise}
Suppose $U$ is a subspace of $V$ and $v_1+U,\dots,v_m+U$ is a basis of $V/U$
and $u_1,\dots,u_n$ is a basis of $U$. Prove that
$v_1,\dots,v_m,u_1,\dots,u_n$ is a basis of $V$.
\end{exercise}
\begin{solution}
First, we show that $v_1,\dots,v_m,u_1,\dots,u_n$ span $V$. Let $v\in V$. Then
$v+U\in V/U$, so we can write
\[ v+U = \lambda_1(v_1+U)+\cdots+\lambda_m(v_m+U) \]
for some scalars $\lambda_1,\dots,\lambda_m$. This gives
$v-(\lambda_1v_1+\cdots+\lambda_mv_m)\in U$. Since $u_1,\dots,u_n$ is a basis
of $U$, we can write $v-(\lambda_1v_1+\cdots+\lambda_mv_m)=\mu_1u_1+\cdots
+\mu_nu_n$ for some scalars $\mu_1,\dots,\mu_n$. Hence,
$v=\lambda_1v_1+\cdots+\lambda_mv_m+\mu_1u_1+\cdots+\mu_nu_n$, showing that
$v_1,\dots,v_m,u_1,\dots,u_n$ span $V$.

Next, we show linear independence. Suppose
$\lambda_1v_1+\cdots+\lambda_mv_m+\mu_1u_1+\cdots+\mu_nu_n=0$. Then
$\lambda_1v_1+\cdots+\lambda_mv_m\in U$. Then
\[ \lambda_1(v_1+U)+\cdots+\lambda_m(v_m+U) = 0+U. \]
Since $v_1+U,\dots,v_m+U$ are linearly independent in $V/U$, we must have
$\lambda_1=\cdots=\lambda_m=0$. Then the original equation reduces to
$\mu_1u_1+\cdots+\mu_nu_n=0$, and since $u_1,\dots,u_n$ are linearly
independent, we have $\mu_1=\cdots=\mu_n=0$. Therefore,
$v_1,\dots,v_m,u_1,\dots,u_n$ are linearly independent, and hence form a basis
of $V$.
\end{solution}

\begin{exercise}
Suppose $\varphi\in\Lin(V,\F)$ and $\varphi\ne0$. Prove that
$\dim V/(\nul\varphi)=1$.
\end{exercise}
\begin{solution}
From \ref{ax:3.107} in the text, we know that $V/\nul\varphi\cong\range\varphi$.
Since $\varphi\ne0$, we have $\range\varphi\ne\{0\}$, and since $\range\varphi$
is a subspace of $\F$, it follows that $\dim\range\varphi=1$. This implies
$\range\varphi=\F$, and hence $\dim V/(\nul\varphi)=\dim\range\varphi=1$.
\end{solution}

\begin{exercise}
Suppose $U$ is a subspace of $V$ such that $\dim V/U=1$. Prove that there exists
$\varphi\in\Lin(V,\F)$ such that $\nul\varphi=U$.
\end{exercise}
\begin{solution}
Since $\dim V/U=\dim\F=1$, there exists some isomorphism
$\psi\in\Lin(V/U,\F)$. Let $\pi\colon V\to V/U$ be the quotient map. Define
$\varphi=\psi\circ\pi\in\Lin(V,\F)$. Suppose $\varphi(v)=0$. Then
$\psi\bigl(\pi(v)\bigr)=0$, which implies $\pi(v)=0$ since $\psi$ is an
isomorphism. Hence, $v\in U$, and we have $\nul\varphi\subseteq U$. Conversely,
if $v\in U$, then $\pi(v)=0$, so $\varphi(v)=\psi\bigl(\pi(v)\bigr)=0$.
Therefore, $\nul\varphi=U$.
\end{solution}

\begin{exercise}
Suppose that $U$ is a subspace of $V$ such that $V/U$ is finite-dimensional.
\begin{parts}
\item Show that if $W$ is a finite-dimensional subspace of $V$ and $V=U+W$, then
  $\dim W\ge\dim V/U$.
\item Prove that there exists a finite-dimensional subspace $W$ of $V$ such that
  $\dim W=\dim V/U$ and $V=U\oplus W$.
\end{parts}
\end{exercise}
\begin{solution}\leavevmode
\begin{parts}
\item Suppose $W$ is a finite-dimensional subspace of $V$ and $V=U+W$. Let
  $\pi\colon V\to V/U$ be the quotient map, and consider the restriction
  $\pi|_W\colon W\to V/U$. To see that $\range(\pi|_W)=V/U$, note that for any
  $v+U\in V/U$, we can write $v=u+w$ for some $u\in U$ and $w\in W$. Then
  $\pi(v)=\pi(u+w)=\pi(u)+\pi(w)=0+\pi(w)=\pi(w)$, so
  $v+U=\pi(w)\in\range(\pi|_W)$. Hence, $\range(\pi|_W)=V/U$. From the
  fundamental theorem of linear maps, we have
  \begin{align*}
  \dim W &= \dim\nul\pi|_W+\dim\range(\pi|_W)\\
         &= \dim\nul\pi|_W+\dim V/U \ge \dim V/U.
  \end{align*}
\item Let $v_1+U,\dots,v_m+U$ be a basis of $V/U$. We claim that
  $v_1,\dots,v_m$ are linearly independent. Suppose
  $\lambda_1v_1+\cdots+\lambda_mv_m=0$ for some scalars
  $\lambda_1,\dots,\lambda_m$. Then
  \[ \lambda_1v_1+\cdots+\lambda_mv_m\in U. \]
  But $\lambda_1(v_1+U)+\cdots+\lambda_m(v_m+U)=0$ in $V/U$, and since
  $v_1+U,\dots,v_m+U$ are linearly independent, it follows that
  $\lambda_1=\cdots=\lambda_m=0$. Hence, $v_1,\dots,v_m$ are linearly
  independent.

  Let $W=\spn\{v_1,\dots,v_m\}$. Then $\dim W=m=\dim V/U$. For any $v\in V$, we
  can write
  \[ v+U = \lambda_1(v_1+U)+\cdots+\lambda_m(v_m+U). \]
  This implies that
  \[ v-(\lambda_1v_1+\cdots+\lambda_mv_m)\in U, \]
  so $v\in U+W$. Hence, $V=U+W$. To see that $U\cap W=\{0\}$, let
  $w\in U\cap W$. Then $w\in W$, so $w=\lambda_1v_1+\cdots+\lambda_mv_m$ for
  some scalars $\lambda_1,\dots,\lambda_m$. But $w\in U$, so $w+U=0$ in $V/U$.
  Hence,
  \[ \lambda_1(v_1+U)+\cdots+\lambda_m(v_m+U) = 0 \]
  in $V/U$. Since $v_1+U,\dots,v_m+U$ are linearly independent, it follows that
  $\lambda_1=\cdots=\lambda_m=0$. Hence, $w=0$, and we have $U\cap W=\{0\}$.
  Therefore, $V=U\oplus W$. \qedhere
\end{parts}
\end{solution}

\begin{exercise}
Suppose $T\in\Lin(V,W)$ and $U$ is a subspace of $V$. Let $\pi$ denote the
quotient map from $V$ onto $V/U$. Prove that there exists $S\in\Lin(V/U,W)$ such
that $T=S\circ\pi$ if and only if $U\subseteq\nul T$.
\end{exercise}
\begin{solution}
Suppose $S\in\Lin(V/U,W)$ is such that $T=S\circ\pi$. Then for any $u\in U$, we
have $T(u)=S\bigl(\pi(u)\bigr)=S(0)=0$, so $U\subseteq\nul T$.

Conversely, suppose $U\subseteq\nul T$. Define $S\colon V/U\to W$ by
$S(v+U)=T(v)$. This is well-defined because if $v+U=v'+U$, then
$v-v'\in U\subseteq\nul T$, so $T(v)=T(v')$. It is straightforward to check that
$S$ is linear and that $T=S\circ\pi$.
\end{solution}
